\chapter{GRAPE: Theory, Gradients, and Implementation}
\label{chap:grape}

One of the prominent techniques for implementing quantum optimal control is Gradient Ascent Pulse Engineering (GRAPE). The central goal of optimal control is to adjust control parameters in such a way that the infidelity of the desired state transfer or gate operation is minimized. In GRAPE, the minimization is based on gradient descent and thus generally requires the evaluation of gradients. Automatic differentiation (AD) \cite{baydin2018automatic} is a convenient tool that has also made an impact on quantum optimal control with GRAPE \cite{pacleung2017speedup,pacqoc,abdelhafez2019gradient,grapead}.
Internally, AD builds a computational graph and applies the chain rule to automatically compute the gradient of a given function. 

However, the convenience of AD comes at the cost of memory required for storing the full computational graph. The use of semi-automatic differentiation (semi-AD) \cite{semiadgoerz2022quantum} partially reduces the memory overhead compared to full automatic differentiation (full-AD). This reduction can still be insufficient to tackle optimal control for quantum systems of rapidly growing size \cite{arute2019,gong2021}. This motivates us to revisit the original GRAPE scheme \cite{nmrkhaneja2005optimal} which is based on hard-coded gradients (HG), and has the smallest possible memory cost. We perform a scaling analysis of memory usage and runtime, and thereby develop a decision tree guiding the optimal choice of strategy among HG, semi-AD and full-AD. We exemplify the scaling behavior with benchmarks for concrete optimization tasks.

The outline of this chaper is as follows. Sec.\ \ref{sec:grape-theory} briefly reviews GRAPE and introduces necessary notation. In Sec.\ \ref{sec:gradient-strategies}, we review three strategies of gradient computation: full automatic differentiation, semi-automatic differentiation and hard-coded gradients. Sec.\ \ref{sec:numerical-gradients}illustrate derivation and memory-efficient implementation of the gradients for the most commonly used cost function contributions. Our numerical implementation of HG further improves the efficiency of state propagation. In Sec.\ \ref{sec:scaling-analysis}, we analyze and compare the scaling of runtime and memory usage scaling for HG, AD and semi-AD. Sec.\ \ref{sec:benchmark-studies}present results from concrete benchmark studies. In Sec.\ \ref{sec:decision-tree}, we formulate a decision tree that facilitates the choice of optimal numerical strategy for GRAPE-based quantum optimal control. 
\section{Gradient Ascent Pulse Engineering: Theory and Notation}
\label{sec:grape-theory}

\subsection{Pulse Parametrization and Time Discretization}
\label{subsec:pulse-param}

\subsection{Cost Functions: State-Transfer Infidelity, Gate Infidelity, and Additional Contributions}
\label{subsec:cost-functions}

\subsection{Pulse Optimization Based on Gradient Information}
\label{subsec:pulse-opt-gradient}

\section{Gradient Computation Strategies}
\label{sec:gradient-strategies}

\subsection{Full Automatic Differentiation}
\label{subsec:full-ad}

\subsection{Semi-Automatic Differentiation}
\label{subsec:semi-ad}

\subsection{Hard-Coded Gradients: Analytical Derivation}
\label{subsec:hard-coded-gradients}

\subsubsection{Gradient of State-Transfer Infidelity}
\label{subsubsec:grad-state-transfer}

\subsubsection{Gradient of Gate Infidelity}
\label{subsubsec:grad-gate-infidelity}

\subsubsection{Gradients of Additional Cost Contributions}
\label{subsubsec:grad-additional-cost}

\section{Numerical Evaluation of Hard-Coded Gradients}
\label{sec:numerical-gradients}

\subsection{State Propagation via the Scaling and Squaring Method}
\label{subsec:scaling-squaring}

\subsubsection{Bounding Numerical Errors}
\label{subsubsec:ss-bounding-errors}

\subsubsection{Choosing the Scaling Factor and Truncation Order}
\label{subsubsec:ss-scaling-truncation}

\subsubsection{Numerical Comparison Between Different Implementations of Scaling and Squaring}
\label{subsubsec:ss-comparison}

\subsection{Propagator Derivative Acting on a State via the Auxiliary Matrix Method}
\label{subsec:auxiliary-matrix}

\section{Scaling Analysis of Runtime and Memory Usage for Three Gradient Strategies}
\label{sec:scaling-analysis}

\subsection{Hard-Coded Gradients}
\label{subsec:scale-hard-coded}

\subsection{Full Automatic Differentiation}
\label{subsec:scale-full-ad}

\subsection{Semi-Automatic Differentiation}
\label{subsec:scale-semi-ad}

\section{Benchmark Studies of Runtime and Memory Usage for the Three Gradient Strategies}
\label{sec:benchmark-studies}

\subsection{State-Transfer Benchmarks}
\label{subsec:state-transfer-benchmarks}

\subsection{Gate-Operation Benchmarks}
\label{subsec:gate-operation-benchmarks}

\section{Decision Tree: A Practical Guide for Selecting Optimization Strategies}
\label{sec:decision-tree}
